\section{Expanded outputs, heights, and Gram and moment certificates}
\label{sec:heights}

The preceding sections compare exact decisions with approximation. This
section asks a different question: once an exact object is known to exist,
how long must its ordinary binary encoding be? The objects are a rational
point of a strictly feasible convex body, a rational optimizer, and the
rational matrices that certify optimality or positivity. These are optimal
order-two moment matrices, optimal polynomial Gram matrices, matrices
exposing the optimal Gram face, and positive definite Gram matrices.

All instances below are rational quartics supplied with a positive definite
rational Hessian Gram on the full basis $(v,X\otimes v)$. Convexity is
therefore certified by short rational data, the Hessian is bounded below by
a positive multiple of the identity on all of $\R^n$, and the lower bounds do
not come from
an unrecognized or nearly degenerate convexity hypothesis. The answers
depend on the representation and on the shape of the certificate:
\begin{enumerate}[label=(\roman*)]
\item a compact convex body with nonempty interior can have only rational
  points with $\Omega(n2^{n/2})$ bits (Theorem~\ref{thm:heights-witness});
\item a unique rational optimizer in $[-1,1]^n$ can require
  $2^{n/2-1}\log_25$ denominator bits, even when the Hessian at the optimizer
  has polynomially bounded condition number
  (Theorem~\ref{thm:rational-height});
\item a strict Hessian Gram makes the order-two moment relaxation exact with
  a unique rank-one optimum, and the optimizer field is then exactly the
  least field for maximal-rank optimal Grams and for nonzero positive
  semidefinite matrices exposing the optimal Gram face
  (Theorem~\ref{thm:heights-moment});
\item consequently, a long rational optimizer forces long entries in every
  maximal-rank rational optimal Gram, in the optimal moment matrix, and in
  every rational exposing matrix, while short optimal Grams of lower rank
  remain (Corollary~\ref{cor:heights-gram});
\item positive definite rational polynomial Grams always exist with
  $\poly(L)2^{O(n)}$ bits for strictly positive certified quartics, and
  $\Omega(n2^{n/2})$ bits are sometimes necessary although a short singular
  Gram exists (Theorem~\ref{thm:interior-gram});
\item for every certified quartic, a polynomial-size shared rational circuit,
  constructed without a sign oracle, gives a strictly feasible point when the
  minimum is negative and a positive definite Gram when the minimum is
  positive; over general certified quartics, validating either output is as
  hard as the exact decision itself
  (Proposition~\ref{prop:upper-circuit-witness} and
  Theorem~\ref{thm:circuit-gram}). The rational optimizers of (ii) are
  computed by short squaring circuits, and short circuits built on these also
  give the optimal moment matrix, the exposing matrices, and a maximal-rank
  optimal Gram. A circuit with rational constants and rational operations has
  rational values, so it cannot describe the irrational optimizers and
  certificates covered by (iii).
\end{enumerate}

These results matter for exact certification in mixed-integer and
polynomial optimization. A solver that must print an exact rational feasible
point, an exact rational optimizer, or a strictly feasible rational Gram
certificate can be forced to print exponentially many bits on small,
perfectly certified convex instances. Approximation does not suffer from this
effect. Shared rational circuits avoid it for strictly feasible points and
positive definite Grams, whenever a certified quartic has them, and for the
long rational optimizers constructed here. For the first two the price is
validation: over general certified quartics, checking the circuit output is
as hard as the exact decision. The certificate shape also matters. Methods
that seek a maximal-rank or interior Gram, or that perform facial reduction while
retaining every real feasible Gram, must reproduce the optimizer's
arithmetic, whereas a lower-rank certificate can be short.

Three limits apply throughout. First, every lower bound concerns
\emph{expanded} rational encodings. None excludes short algebraic,
circuit, or other implicit descriptions, and none is a lower bound on the
complexity of a decision problem or of approximation. Second, the
constructed families have total input length $L_k$ bounded by a polynomial
in the family parameter $k$, and their dimension is $n=\Theta(k)$. The lower
bounds are exponential in $n$ and hence superpolynomial in $L_k$, but they
are not of the form $2^{\Omega(L_k)}$. Third, only the stated certificate
shapes are covered. In particular, the lower bounds for Gram matrices concern
maximal-rank, exposing, or positive definite Grams; the same instances have
short rational Gram matrices of lower rank.

\subsection{Encodings, heights, and Taylor Gram matrices}
\label{sec:heights-conventions}

A rational number $r$ is written in lowest terms as $a/b$ with $b\ge1$, and
$\operatorname{den}(r)=b$. Its height is $\height(r)=\log_2\max\{|a|,b\}$;
for $r=0$ put $\height(r)=0$. A rational vector or matrix has \emph{entry
height at most $B$} if every entry has height at most $B$. Any ordinary
binary encoding of a rational vector uses at least
$\sum_i\log_2\operatorname{den}(r_i)$ bits, so lower bounds on denominators
are lower bounds on expanded output length. A shared rational circuit is a
directed acyclic graph of rational arithmetic gates, as in
Section~\ref{sec:models}; its size counts gates and the binary length of its
constants. A circuit of size $O(k)$ can compute a rational whose reduced
fraction has more than $2^k$ bits, for instance $2^{-2^k}$ by repeated
squaring.

For $X\in\R^n$ let $z(X)$ be the vector of all monomials of degree at most two,
with the constant monomial first; its length is $N=\binom{n+2}{2}$. A
symmetric matrix $Q$ is a (polynomial) Gram matrix of $f$ if
$f=z^{\mathsf T}Qz$. Put $w(X,v)=(v,X\otimes v)\in\R^{n+n^2}$. A symmetric
matrix $A$ of order $n+n^2$ is a \emph{full Hessian Gram} of a polynomial $f$
of degree at most four if
\begin{equation}
 v^{\mathsf T}\nabla^2f(X)v=w(X,v)^{\mathsf T}A\,w(X,v)
 \qquad(X,v\in\R^n).
 \label{eq:heights-hessian-gram}
\end{equation}
A \emph{certified quartic} is a pair $(f,A)$ with $f\in\Q[X]$ of degree at
most four and $A$ a rational full Hessian Gram with $A\succ0$. Validity is
checked in polynomial time by coefficient comparison and rational
elimination. If $A\succ0$ has order $s=n+n^2$, then the number
$\mu_A=\det A/(\tr A)^{s-1}$, denoted $\rho(A)$ in
Lemma~\ref{lem:models-det-trace}, is a positive rational of polynomial bit
length with $A\succeq\mu_AI$, and
$\nabla^2f\succeq\mu_AI$ on $\R^n$
(Lemma~\ref{lem:models-hessian-gram}(a)). Thus $f$ is coercive and has a unique
minimizer, denoted $p$. Polynomial Gram matrices and Hessian Gram matrices
are different objects: the former represent $f$ in the basis $z$, the latter
represent its Hessian biform in the basis $w$.

The following integration lemma is the common source of all exact Gram
matrices used in this paper. For $a\in\R^n$ write $d=X-a$ and define matrices
$C_U(a),C_V(a)$ with $n+n^2$ rows and $N$ columns by
\[
 (d,\,a\otimes d)=C_U(a)\,z(X),\qquad (0,\,d\otimes d)=C_V(a)\,z(X).
\]
Their entries are integer polynomials of degree at most two in the
coordinates of $a$. For a symmetric matrix $B$ of order $n+n^2$ define the
\emph{Taylor Gram}
\begin{equation}
 \mathcal T_a(B)=\tfrac12\bigl(C_U(a)+\tfrac13C_V(a)\bigr)^{\mathsf T}B
          \bigl(C_U(a)+\tfrac13C_V(a)\bigr)
   +\tfrac1{36}C_V(a)^{\mathsf T}B\,C_V(a).
 \label{eq:heights-taylor-gram}
\end{equation}

\begin{lemma}[Taylor integration of a full Hessian Gram]
\label{lem:taylor-sos}
Let $f\in\R[X_1,\ldots,X_n]$ have degree at most four and a real full Hessian
Gram $A$. Fix $a\in\R^n$, and put $d=X-a$ and $b=\nabla f(a)$.
\begin{enumerate}[label=(\alph*)]
\item \emph{Identity.}
  $f(X)=f(a)+b^{\mathsf T}d+z(X)^{\mathsf T}\mathcal T_a(A)z(X)$ identically.
  The identity rests on
  \[
   \int_0^1(1-t)\,(U+tV)^{\mathsf T}A\,(U+tV)\,dt
   =\tfrac12(U+V/3)^{\mathsf T}A(U+V/3)+\tfrac1{36}V^{\mathsf T}AV
  \]
  for $U=(d,a\otimes d)$ and $V=(0,d\otimes d)$, whose scalar form is
  $\int_0^1(1-t)(u+tv)^2dt=\frac12(u+v/3)^2+(v/6)^2$.
\item \emph{Positivity and the strict kernel.} If $A\succeq0$ then
  $\mathcal T_a(A)\succeq0$. If $A\succ0$ then
  $\ker\mathcal T_a(A)=\R z(a)$; equivalently, after the invertible change to
  the translated basis $(1,(d^\alpha)_{1\le|\alpha|\le2})$, the Taylor Gram is
  $0\oplus J$ with $J$ positive definite on the nonconstant monomials.
\item \emph{Gradient completion.} Suppose $A\succeq\mu I$ with $\mu>0$. Put
  $\bar A=A-\mu\diag(I_n,0)$, $c_a=f(a)-\norm{b}^2/(2\mu)$, let
  $C_{\rm aff}(a)$ be the coefficient matrix with
  $d+b/\mu=C_{\rm aff}(a)z(X)$, and let $e_0$ select the constant monomial.
  Then $f=z^{\mathsf T}\mathcal Q_a z$ for
  \begin{equation}
   \mathcal Q_a=\mathcal T_a(\bar A)
     +\tfrac\mu2C_{\rm aff}(a)^{\mathsf T}C_{\rm aff}(a)
     +c_a\,e_0e_0^{\mathsf T}.
   \label{eq:heights-completed-gram}
  \end{equation}
  If $c_a\ge0$ then $\mathcal Q_a\succeq0$, and if $c_a>0$ then
  $\mathcal Q_a\succ0$.
\item \emph{Fields.} Let $E\subseteq\R$ be a field containing the
  coefficients of $f$, the entries of $A$, and the coordinates of $a$ (and
  $\mu$ in part (c)). Then $C_U(a)$, $C_V(a)$, $C_{\rm aff}(a)$,
  $\mathcal T_a(A)$, and $\mathcal Q_a$ have entries in $E$. In particular,
  at a critical point $a$ of $f$ with $A\succeq0$, the matrix
  $\mathcal T_a(A)$ is a positive semidefinite Gram of $f-f(a)$ with entries
  in $E$.
\item \emph{Squares.} If $A=\sum_ks_k\ell_k\ell_k^{\mathsf T}$ with $s_k\ge0$,
  then, with $U,V$ as in (a),
  \[
   f-f(a)-b^{\mathsf T}d=\sum_ks_k\Bigl[2\Bigl(\frac{\ell_k^{\mathsf T}(U+V/3)}2
     \Bigr)^2+\Bigl(\frac{\ell_k^{\mathsf T}V}6\Bigr)^2\Bigr].
  \]
  If $A\succeq0$ is rational, the $s_k$ and $\ell_k$ can be chosen rational
  with each $s_k$ a sum of rational squares, so $f-f(a)-b^{\mathsf T}d$ is a
  sum of squares of quadratics whose coefficients are rational polynomials of
  degree at most two in the coordinates of $a$. In
  particular, at a critical point $a$, $f-f(a)$ is a sum of squares over every
  subfield $E\subseteq\R$ that contains the coordinates of $a$.
\end{enumerate}
\end{lemma}

The proof is in Appendix~\ref{app:heights-taylor}. The integral identity is
the classical argument that an SOS-convex polynomial minus its first-order
Taylor expansion is a sum of squares \cite{HeltonNie2010,Lasserre2009}. The
lemma adds what the exact constructions need: an explicit Gram matrix, its
exact kernel in the strict case, gradient completion with an explicit
constant, and field bookkeeping.

\begin{remark}[Gram fields and square factors]
\label{rem:heights-gram-field}
Part (d) concerns the field of \emph{Gram entries}, and part (e) the field of
\emph{square factors}; the two need separate arguments. Over $\Q$ a positive
semidefinite rational Gram yields a sum of squares of rational polynomials:
rational $LDL^{\mathsf T}$ elimination gives nonnegative rational weights, and
each positive rational weight is a sum of rational squares
(Lemma~\ref{lem:models-rational-squares}). This is why part (e) needs a
rational Hessian Gram but no rationality of $a$. Over a general real field
$E$ the same elimination gives a sum of squares with nonnegative weights in
$E$, but an unweighted sum of squares with coefficients in $E$ need not
exist. For
example, with $E=\Q(\sqrt2)\subset\R$, the polynomial $\sqrt2\,X_1^2$ has a
positive semidefinite Gram over $E$, namely the matrix whose only nonzero
entry is the diagonal entry $\sqrt2$ at the monomial $X_1$, but an identity
$\sqrt2\,X_1^2=\sum_iq_i^2$ with $q_i\in E[X_1]$ is impossible: comparing the
coefficients of $X_1^2$ and applying the conjugation $\sqrt2\mapsto-\sqrt2$
would make $-\sqrt2$ a sum of squares of real numbers. Field statements in this
section are therefore always about Gram entries; least fields of square
factorizations are treated in Section~\ref{sec:fields}.
\end{remark}

\subsection{Strictly feasible bodies with only long rational points}
\label{sec:heights-witness}

A strictly feasible convex body always contains rational points, because it
has an interior. The next theorem shows that it need not contain a short
one, even when it is the zero sublevel set of a single strongly convex
quartic with a short rational convexity certificate and lies in a fixed box.

\paragraph{Construction.}
Let $k\ge1$, $n=2k$, and $M_k=1000^{k+3}$. A doubly exponentially small
positive number is generated by $k$ cube roots:
\begin{equation}
 \delta_0=M_k^{-1},\qquad
 \delta_i=(1+3\delta_{i-1}^2)^{1/3}-1\quad(1\le i\le k).
 \label{eq:heights-tiny-chain}
\end{equation}
Since $(1+t)^3\ge1+3t$ for $t\ge0$, induction gives
$0<\delta_i\le\delta_{i-1}^2$, hence $0<\delta_k\le M_k^{-2^k}$. With
$\xi_i=1+\delta_i$ the gate equations are
\begin{equation}
 \xi_1^3=1+3M_k^{-2},\qquad
 \xi_i^3=3\xi_{i-1}^2-6\xi_{i-1}+4\quad(2\le i\le k).
 \label{eq:heights-cube-gates}
\end{equation}
Write the coordinates as $X=(x_1,y_1,\ldots,x_k,y_k)$ and put
$p=(\xi_1,\xi_1^2,\ldots,\xi_k,\xi_k^2)$. Lemma~\ref{lem:heights-cube-chain}
in the appendix constructs, in time polynomial in $k$, a rational quartic
$F_k$ given as a sum of $n+1$ squares of rational quadratics, with
$\nabla^2F_k\succeq\frac32I$, unique zero $p$, and a rational positive
definite full Hessian Gram $A_k$, all of polynomial bit length. The number
$\delta_k$ is never printed; it is the value at $p$ of the rational affine
function $u_k(X)=x_k-1$. Let $\iota_k\in\R^{n+n^2}$ be the coordinate vector
of $v_{x_k}$ in the first block of $w(X,v)$, and put
\begin{equation}
 \lambda_k=\left\lceil\frac3{\mu_{A_k}}\right\rceil,\qquad
 g_k=\lambda_kF_k-u_k^2,\qquad
 B_k=\lambda_kA_k-2\iota_k\iota_k^{\mathsf T}.
 \label{eq:heights-witness-family}
\end{equation}
The Hessian biform of $u_k^2$ is $2v_{x_k}^2$, so $B_k$ is a full Hessian Gram
of $g_k$, and $B_k\succeq3I-2I=I$. The integer $\lambda_k$ can be
exponentially large, but its binary length is polynomial in $k$.

\begin{theorem}[Strict feasibility without short rational points]
\label{thm:heights-witness}
For every $k\ge1$ the certified quartic $(g_k,B_k)$ of
\eqref{eq:heights-witness-family} in $n=2k$ variables is computable in time
polynomial in $k$, and its total binary length $L_k$ is bounded by a
polynomial in $k$. Let $S_k=\{X\in\R^n:g_k(X)\le0\}$.
\begin{enumerate}[label=(\alph*)]
\item $\nabla^2g_k\succeq I$ on $\R^n$. The set $S_k$ is compact and strictly
  convex, has nonempty interior, and is contained in $(0.99,1.01)^n$.
\item Every rational point $X\in S_k$, including boundary points, satisfies
  the following. If $X_1=a/b$ with integers $a$ and $b\ge1$, then
  \[
   \log_2b>\tfrac13\bigl((2^k-2)(k+3)\log_21000-4\bigr).
  \]
  In particular $\log_2b>k2^k=\frac n2\,2^{n/2}$ for every $k\ge2$, so
  every rational point of $S_k$ has ordinary binary length greater than
  $\frac n2\,2^{n/2}$.
\item A shared rational circuit of size polynomial in $k$ for a point
  $\hat x\in\Q^n$ with $g_k(\hat x)<0$ is computable in time polynomial in
  $k$, without sign tests (Proposition~\ref{prop:upper-circuit-witness}).
\end{enumerate}
\end{theorem}

The proof is in Appendix~\ref{app:heights-witness}. Since $F_k(p)=0$ and
$\nabla F_k(p)=0$, the point $p$ has $g_k(p)=-\delta_k^2$ and
$\norm{\nabla g_k(p)}=2\delta_k$. Strong convexity then confines every point
of $S_k$, including boundary points, to the ball of radius
$(2+\sqrt6)\delta_k<5M_k^{-2^k}$ about $p$. The first coordinate of $p$ is the
irrational number $\xi_1$, a root of the integer cubic
$M_k^2T^3-(M_k^2+3)$. For a rational point with first coordinate $a/b$, the
integer $M_k^2a^3-(M_k^2+3)b^3$ is nonzero, so its absolute value is at least
one; the difference-of-cubes factorization bounds it by
$16b^3M_k^{2-2^k}$, which forces the denominator bound.

The argument needs no Diophantine approximation theorem: the irrational
coordinate satisfies an explicit integer cubic with small coefficients, and a
nonzero integer has absolute value at least one. The obstruction is precision
near that coordinate, not a large coordinate magnitude, empty interior, or a
degenerate convexity certificate. The body lies in a ball of radius below
$5M_k^{-2^k}$ about $p$, so the theorem is also an explicit example in which
strict feasibility supplies no quantitative Slater radius of polynomial bit
length. The localization applies to every point of $S_k$, so the bound is not
a statement about one particular rounded point.

Large rational witnesses for convex systems are classical. Khachiyan's
repeated-squaring system $x_i\ge x_{i+1}^2$, $x_m\ge2$, recalled by Pataki
and Touzov \cite{PatakiTouzov2024}, is a strictly feasible convex quadratic
system whose rational solutions have exponentially many bits, because its
coordinates are huge. Theorem~\ref{thm:heights-witness} adds the
restrictions of a single quartic inequality with a supplied strict rational
full Hessian Gram, a compact full-dimensional feasible set, and bounded
coordinates. In the other direction, a nonempty set defined by one strict
integer polynomial inequality of degree $d$ with coefficients of bit length
$\tau$ contains a rational point with coordinates of bit length
$\tau d^{O(n)}$ \cite[Theorem~4.1.2]{BasuPollackRoy1996}; see also
\cite[Proposition~2.5]{SafeyElDinZhi2010}. Applied to $\{g_k<0\}$ this gives
rational points with $\poly(k)2^{O(k)}$ bits, so the exponential dependence
on the dimension in Theorem~\ref{thm:heights-witness} is necessary and
attained up to the constants in the exponent. The irrational singleton
examples of Section~\ref{sec:algebraic} have no rational feasible point at
all; here rational points exist but are long.

\subsection{Rational optimizers with exponentially long denominators}
\label{sec:heights-rational}

A rational optimizer can also be long. The construction powers a rational
point of the unit circle. Repeated squaring of $(3+4\mathrm i)/5$ never
cancels a factor of five, so the denominators grow doubly exponentially
while the coordinates stay in $[-1,1]$.

\begin{lemma}[Denominators of the circle chain]
\label{lem:heights-circle-denominators}
For an integer $m\ge1$ write $(3+4\mathrm i)^m=A_m+\mathrm iB_m$ with
$A_m,B_m\in\Z$. Then $A_m\equiv3$ and $B_m\equiv4\pmod5$. Consequently the
real and imaginary parts of $((3+4\mathrm i)/5)^m$ have reduced denominators
exactly $5^m$, and their absolute values are at most one.
\end{lemma}

\begin{proof}
In $\Z[\mathrm i]$ we have $(3+4\mathrm i)^2=-7+24\mathrm i\equiv3+4\mathrm i$
modulo $5\Z[\mathrm i]$. By induction $(3+4\mathrm i)^m\equiv3+4\mathrm i$
for all $m\ge1$. Since $5$ divides neither $3$ nor $4$, the fractions
$A_m/5^m$ and $B_m/5^m$ are in lowest terms. The modulus of
$((3+4\mathrm i)/5)^m$ is one.
\end{proof}

Let $k\ge0$ and $n=2(k+1)$, and write the coordinates as
$X=(x_0,y_0,\ldots,x_k,y_k)$. Put $z_j=((3+4\mathrm i)/5)^{2^j}=a_j+\mathrm ib_j$.
The rational quadratic residuals
\begin{equation}
 r_0=x_0-\tfrac35,\quad s_0=y_0-\tfrac45,\quad
 r_j=x_j-x_{j-1}^2+y_{j-1}^2,\quad s_j=y_j-2x_{j-1}y_{j-1}\quad(1\le j\le k)
 \label{eq:heights-circle-residuals}
\end{equation}
have the unique common real zero $p=(a_j,b_j)_{j=0}^k$; they express complex
squaring with coefficients of constant size. Every circle relation
$x_j^2+y_j^2-1$ also vanishes at $p$, and the circle relations supply the
positive curvature that the residuals lack. Lemma~\ref{lem:heights-circle-chain}
in the appendix combines these quadratics, with coefficients approximated
only where approximation preserves the exact zero, into a quartic $f_k$ via
Lemma~\ref{lem:quartic-realization}. A rescaled version with circles of radius
$2^{-j}$ controls the Hessian at the optimizer.

\begin{theorem}[Bounded rational optimizers with long denominators]
\label{thm:rational-height}
For every integer $k\ge0$ put $n=2(k+1)$. In time polynomial in $k$ one can
construct rational quartics $f_k$ and $\tilde f_k$ in $n$ variables, each
given as a sum of $n+1$ squares of rational quadratics and each supplied with
a rational positive definite full Hessian Gram, all of total binary length
bounded by a polynomial in $k$, with the following properties.
\begin{enumerate}[label=(\alph*)]
\item $\nabla^2f_k\succeq\frac32I$ and $\nabla^2\tilde f_k\succeq\frac32I$ on
  $\R^n$. Both minima are zero and are attained only at points $p$ and
  $\tilde p$, respectively.
\item $p\in\Q^n\cap[-1,1]^n$, and the last two coordinates of $p$ have reduced
  denominator exactly $5^{2^k}$.
\item $\tilde p\in\Q^n$, $\norm{\tilde p}^2<\frac43$, the last two coordinates
  of $\tilde p$ have reduced denominators divisible by $5^{2^k}$, and
  \[
   2I\preceq\nabla^2\tilde f_k(\tilde p)\preceq\bigl(8(k+1)^2+1\bigr)I.
  \]
  Thus the condition number of the Hessian at the optimizer is at most
  $4(k+1)^2+\frac12$.
\end{enumerate}
Consequently, writing either optimizer as ordinary reduced fractions requires
at least $2^k\log_25=2^{n/2-1}\log_25$ bits in one coordinate. This is
exponential in $n$ and superpolynomial in the input length, but it is not of
the form $2^{\Omega(L_k)}$.
\end{theorem}

The proof is in Appendix~\ref{app:heights-circle}. Both optimizers have
short exact descriptions: the shared circuit computing
$((3+4\mathrm i)/5)^{2^k}$ by $k$ squarings, followed by a rescaling for
$\tilde p$. They can also be approximated to any requested precision in
polynomial time. The theorem therefore separates the size of an exact
expanded rational output from both exact implicit output and approximation.
It is also different from the rational-optimizer hardness of
Theorem~\ref{thm:rational-optimizer}: that result concerns the cost of
\emph{comparing} a rational optimizer coordinate with zero, under a
rationality promise, whereas Theorem~\ref{thm:rational-height} is an
unconditional bound on the \emph{length} of the optimizer itself.

Part (c) shows that a long rational optimizer is compatible with good local
conditioning: the Hessian at $\tilde p$ has condition number $O(n^2)$. This is
a statement about the Hessian at the optimizer only. The coefficients of
$\tilde f_k$ have magnitudes up to $2^{O(k)}$, and no uniform smoothness
bound on a fixed neighborhood of $\tilde p$, no bound on higher derivatives,
and no bound on the conditioning of rational reconstruction is asserted.
Multiplying $f_k$ by a sufficiently large integer square normalizes its
Hessian Gram to be at least the identity without changing $p$.

The norm and approximation guarantees for convex polynomial optimization of
Slot, Steurer, and Wiedmer \cite{SlotSteurerWiedmer2025} are compatible
with Theorem~\ref{thm:rational-height}: they bound magnitudes and
approximation cost, not denominators. Jiang's exact algorithm for convex
functions with rational minimizers \cite{Jiang2021} is parameterized by a
bound on the common denominator of the minimizer; for $f_k$ the least common
denominator of the coordinates of $p$ is $5^{2^k}$, although the coordinates
are bounded by one. Exponentially
long rational local minimizers of nonconvex cubics, caused by large
coordinate magnitudes, appear in Zhang's thesis \cite[Example~2.5.3]{Zhang2020}.
We do not claim that powering a rational point of the circle is new; the
contribution is its realization as the unique optimizer of a globally
strongly convex quartic with short rational square factors and a short
strict Hessian certificate.

\subsection{Exact moment relaxations and the arithmetic of optimal Grams}
\label{sec:heights-moment}

For a certified quartic $f$ consider the unconstrained order-two moment
relaxation and its sum-of-squares dual. For a sequence
$y=(y_\alpha)_{|\alpha|\le4}$ let $\Lambda_y(f)=\sum_\alpha f_\alpha y_\alpha$
and $M_2(y)=(y_{\alpha+\beta})_{|\alpha|,|\beta|\le2}$, indexed like $z$. The
programs are
\begin{equation}
 \text{(P)}\ \ \min\{\Lambda_y(f):y_0=1,\ M_2(y)\succeq0\},
 \qquad
 \text{(S)}\ \ \max\{\gamma:f-\gamma=z^{\mathsf T}Qz,\ Q\succeq0\}.
 \label{eq:heights-moment-programs}
\end{equation}
Their data are rational. Let $p$ be the minimizer of $f$, $f_*=f(p)$,
$e=z(p)$, and let $K=\Q(p_1,\ldots,p_n)\subseteq\R$ be the field generated by
the coordinates of $p$. An \emph{optimal Gram} is a positive semidefinite
Gram of $f-f_*$; the set of all of them is denoted $\mathcal G_*$, and the
\emph{optimal-level Gram system} is $f-f_*=z^{\mathsf T}Qz$, $Q\succeq0$. A
nonzero positive semidefinite $Z$ \emph{exposes} $\mathcal G_*$ if
$\ip ZQ=\tr(ZQ)=0$ for every $Q\in\mathcal G_*$. A facial-reduction step
\cite{BorweinWolkowicz1981} restricts a semidefinite system to the face cut
out by such a matrix of the form $\mathcal A^*(y)$, where $\mathcal A$ is the
coefficient map of the system and $y$ annihilates its right-hand side; the
\emph{singularity degree} is the least number of steps after which the
restricted system has a feasible point in the relative interior of its
face.

\begin{theorem}[Moments, maximal-rank Grams, and exposing matrices]
\label{thm:heights-moment}
Let $(f,A)$ be a certified quartic in $n$ variables. Put
$r=\dim_\Q\Span_\Q\{p^\alpha:|\alpha|\le2\}$.
\begin{enumerate}[label=(\alph*)]
\item \emph{Maximal rank.} The Taylor Gram $Q_*=\mathcal T_p(A)$ is an optimal
  Gram with entries in $K$ and $\ker Q_*=\R e$. Every optimal Gram annihilates
  $e$, so $N-1$ is the largest rank of an optimal Gram.
\item \emph{Exactness and uniqueness.} Both programs in
  \eqref{eq:heights-moment-programs} attain the value $f_*$. The only optimal
  solution of (P) is $y_\alpha=p^\alpha$, that is, $M_2(y)=ee^{\mathsf T}$.
  The optimal pair $(ee^{\mathsf T},Q_*)$ is strictly complementary, and both
  programs have strictly feasible points.
\item \emph{Least field of maximal-rank Grams.} For a field
  $E\subseteq\R$, some optimal Gram of rank $N-1$ has all entries in $E$ if
  and only if $K\subseteq E$.
\item \emph{Rational optimal Grams.} Every rational optimal Gram has rank at
  most $N-r$. If $p\notin\Q^n$ then $r\ge2$, so no rational optimal Gram is
  strictly complementary to the optimal moment matrix.
\item \emph{Faces and exposing matrices.} The smallest face of the positive
  semidefinite cone containing $\mathcal G_*$ is
  $\mathcal F_p=\{Q\succeq0:Qe=0\}$, and $Q_*$ lies in its relative interior.
  A nonzero positive semidefinite $Z$ exposes $\mathcal G_*$ if and only if
  $Z=c\,ee^{\mathsf T}$ with $c>0$. The entries of every such $Z$ generate a
  field containing $K$, and $ee^{\mathsf T}$ has entries in $K$. The matrix
  $ee^{\mathsf T}$ is a facial-reduction certificate for the optimal-level
  Gram system, one reduction step reaches $\mathcal F_p$, and the real
  singularity degree of that system is exactly one. If $p\notin\Q^n$, no
  nonzero rational positive semidefinite matrix exposes $\mathcal G_*$.
\end{enumerate}
\end{theorem}

The proof is in Appendix~\ref{app:heights-moment}. Its core is short. The
Taylor Gram of Lemma~\ref{lem:taylor-sos} at the stationary point $p$ has
kernel exactly $\R e$. If $y$ is optimal for (P), then
$\tr(Q_*M_2(y))=\Lambda_y(f-f_*)=0$; two positive semidefinite matrices with
zero trace product have orthogonal ranges, so $M_2(y)$ is a multiple of
$ee^{\mathsf T}$, and $y_0=1$ fixes the multiple. Every moment of degree at
most four is an entry of $M_2(y)$. A maximal-rank Gram over $E$ has a
one-dimensional kernel over $E$, and normalizing its constant coordinate
recovers $e$, hence $p$. The same kernel argument identifies every exposing
matrix.

Exactness of low-order moment relaxations for SOS-convex problems, with
extraction of the optimizer from first-order moments, is due to Lasserre
\cite[Theorems~2.6 and~3.3]{Lasserre2009}. Theorem~\ref{thm:heights-moment}
uses the stronger strict Hessian certificate to determine \emph{all} moments
and the complete kernel of the optimal Gram face. The use of real zeros as
Gram kernel vectors, and of their conjugates to obtain rational kernel
relations, is established in Laplagne's facial-reduction method for exact
SOS decompositions \cite[Proposition~3.2, Section~3.2]{Laplagne2020}; part
(d) is a dimension count for that mechanism. Algebraic, possibly irrational,
exact solutions of degenerate semidefinite programs and their certification
are also established \cite{KolmogorovNaldiZapata2024}. The point of
Theorem~\ref{thm:heights-moment} is the precise identification of the field
forced on each certificate shape, in a setting where the relaxation is exact,
both programs are strictly feasible, strict complementarity holds over the
reals, and the optimal Gram face is reached in one real facial-reduction
step. Field statements concern Gram and moment entries
(Remark~\ref{rem:heights-gram-field}).

Part (e) concerns matrices orthogonal to \emph{all real} optimal Grams. It
does not prevent a procedure that deliberately restricts attention to
rational Gram candidates. Such a procedure may impose the rational kernel
relations of part (d) and work in a smaller face that discards real optimal
Grams; that is a different feasibility problem. The interpretation of
$ee^{\mathsf T}$ as a facial-reduction certificate for a \emph{rational}
optimal-level system needs $f_*\in\Q$. If $f_*\notin\Q$, rational optimal
Grams do not exist at all, because the constant coefficient of $f-f_*$ is
irrational.

\begin{remark}[The strict Hessian hypothesis]
\label{rem:heights-moment-hypothesis}
Strong convexity together with SOS-convexity does not suffice for
part (b). The polynomial $f=X_1^2+X_2^2+X_2^4$ has
$\nabla^2f=\diag(2,2+12X_2^2)\succeq2I$ and a diagonal SOS Hessian, with
minimizer $0$ and $f_*=0$. The sequence with $y_0=1$, $y_{(4,0)}=t\ge0$, and
all other moments zero has $M_2(y)=\diag(1,0,0,t,0,0)$ in the basis
$(1,X_1,X_2,X_1^2,X_1X_2,X_2^2)$ and objective value zero, so (P) has
infinitely many optimal solutions. This $f$ has no positive definite full
Hessian Gram: its Hessian biform has no $X_1^2v_1^2$ term, so the diagonal
entry of a full Hessian Gram at the coordinate $X_1v_1$ must vanish. The
hypothesis is sufficient but not necessary. The proof uses it only to obtain
an optimal Gram with kernel exactly $\R e$; for instance
$X_1^2+X_2^2+X_1^4+X_2^4$ has no positive definite full Hessian Gram, by the
same coefficient argument, but has the optimal Gram
$0\oplus I_2\oplus\left(\begin{smallmatrix}1&0&-1/2\\0&1&0\\-1/2&0&1
\end{smallmatrix}\right)$, whose kernel is $\R e$.
\end{remark}

\begin{remark}[Large fields]
\label{rem:heights-cyclic-fields}
For the cyclic quartics of Theorem~\ref{thm:algebraic-cyclic}, which are
certified quartics with minimum zero and optimizer field of degree
$d_n=(2^{n+1}-(-1)^{n+1})/3$, parts (c) and (e) show that every maximal-rank
optimal Gram and every exposing matrix of the optimal Gram face has entries
generating a field of degree at least $d_n$ over $\Q$. This is a statement
about fields; it does not exclude short circuit or other implicit
descriptions of these matrices.
\end{remark}

\begin{remark}[Nullvectors of one feasible matrix]
\label{rem:heights-nullvector}
Facial reduction preserves every real feasible point when it uses a vector
in the common kernel of the feasible set. A maximum-rank feasible matrix has
exactly this common kernel: if $G_0,G\succeq0$ then
$\ker(G_0+G)=\ker G_0\cap\ker G$, so a feasible $G$ that did not vanish on
$\ker G_0$ would give the feasible midpoint $(G_0+G)/2$ a larger rank. This is
why part (a) matters. A nullvector of an arbitrary feasible matrix does not
suffice, even when it is rational. Proposition~\ref{prop:bnd-nullvector}
constructs a rational affine semidefinite family with a rational positive
definite member and a singular feasible member with a rational nullvector.
Restricting the family to matrices that annihilate that vector leaves only a
single irrational feasible matrix, removing every rational feasible point.
The construction and full proof are given in Appendix~\ref{app:bnd-nullvector}.
\end{remark}

\subsection{Heights forced by a long rational optimizer}
\label{sec:heights-gram-heights}

When the optimizer is rational, Theorem~\ref{thm:heights-moment} gives
rational maximal-rank Grams and rational exposing matrices, but their
entries must encode the optimizer. For the circle family this makes them
long, although a short rational optimal Gram of smaller rank is part of the
input.

\begin{corollary}[Long maximal-rank, moment, and exposing certificates]
\label{cor:heights-gram}
Let $f_k$ be the quartic of Theorem~\ref{thm:rational-height} in
$n=2(k+1)$ variables, with optimizer $p$ and minimum zero, and let
$N=\binom{n+2}{2}$ and $H_k=2^k\log_25$.
\begin{enumerate}[label=(\alph*)]
\item The unique optimal moment matrix $ee^{\mathsf T}$ has an entry with
  reduced denominator $5^{2^k}$.
\item Rational optimal Grams of rank $N-1$ exist. If such a Gram has entry
  height at most $B$, then
  \[
   B\ge\frac{H_k-\log_2((N-1)!)}{N^2-1}=\Omega(2^k/n^4),
  \]
  and the sum of the binary lengths of all its entries is at least
  $H_k-\log_2((N-1)!)$.
\item If a nonzero rational positive semidefinite matrix exposing the optimal
  Gram face has entry height at most $B$, then $B\ge H_k/2$.
\item $f_k$ has a rational optimal Gram of rank at most $n+1<N-1$ and binary
  length polynomial in $k$, namely the Gram of its supplied square factors.
\end{enumerate}
The same statements hold for $\tilde f_k$, with ``reduced denominator
$5^{2^k}$'' replaced by ``reduced denominator divisible by $5^{2^k}$''.
\end{corollary}

The proof is in Appendix~\ref{app:heights-moment}. For (b), the principal
submatrix of a maximal-rank optimal Gram on the nonconstant monomials is
nonsingular and determines the remaining coordinates of $e=(1,w)$ by Cramer's
rule; a terminal coordinate of $p$ appears in $w$, so its denominator divides
a determinant bounded in terms of the entry heights. For (c), every exposing
matrix is $c\,ee^{\mathsf T}$, and a terminal coordinate of $p$ is the ratio
of two of its entries. These lower bounds are exponential in $n$ and
superpolynomial in the input length; part (d) shows that they cannot be
extended to all optimal Grams.

The phenomenon of short low-rank feasible points coexisting with long
maximal-rank feasible points is classical in semidefinite programming. In
the homogenized Khachiyan system with blocks
$\left(\begin{smallmatrix}x_1&2t\\2t&t\end{smallmatrix}\right)$ and
$\left(\begin{smallmatrix}x_j&x_{j-1}\\x_{j-1}&t\end{smallmatrix}\right)$
for $2\le j\le k$ (compare \cite{PatakiTouzov2024}), the point $t=x_1=\cdots=
x_{k-1}=0$, $x_k=1$ is a short feasible point of rank one. Every feasible
point of full rank has $t>0$, $x_1/t>4$, and $x_j/t>(x_{j-1}/t)^2$, hence
$x_k/t>2^{2^k}$; if $x_k$ and $t$ have height at most $B$ then
$2^{2B}\ge x_k/t$, so $B>2^{k-1}$. Corollary~\ref{cor:heights-gram} does not
claim this general distinction as new. Its content is the realization of the
distinction by the optimal Gram face of a short certified convex quartic with
a bounded rational optimizer, where it concerns maximal-rank, moment, and
exposing certificates simultaneously.

\subsection{Interior Gram matrices}
\label{sec:heights-interior}

A polynomial has a positive definite Gram matrix only if it is strictly
positive, because $z$ contains the constant monomial. For certified quartics
the converse holds: if $\min f>0$, then by continuity some rational center
$a$ satisfies $c_a>0$ in Lemma~\ref{lem:taylor-sos}(c), and $\mathcal Q_a$ is
a rational positive definite Gram. Interior rational Grams are the
certificates produced by methods that round an approximate interior solution
and project it onto the rational coefficient equations
\cite{PeyrlParrilo2008}. The question is how long they must be.

Let $\norm{f}_1$ denote the sum of the absolute values of the coefficients of
$f$. Continuing the construction of Section~\ref{sec:heights-witness}, put
\begin{equation}
 h_k=\lambda_kF_k+u_k^2,\qquad
 \lambda_kA_k+2\iota_k\iota_k^{\mathsf T}\succeq3I,
 \label{eq:heights-interior-family}
\end{equation}
which differs from $g_k$ only in the sign of the square.

\begin{theorem}[Interior Gram matrices]
\label{thm:interior-gram}
\leavevmode
\begin{enumerate}[label=(\alph*)]
\item \emph{Upper bound.} Let $(f,A)$ be a certified quartic in $n$
  variables with $\min f>0$, and let $L$ be the total binary length of $f$
  and $A$. Then $f$ has a rational positive definite Gram matrix of total
  binary length at most $\poly(L)\,2^{O(n)}$, and $f$ is a sum of squares of
  rational quadratics of the same total length.
\item \emph{A strictly positive family.} For every $k\ge1$ the certified
  quartic $(h_k,\lambda_kA_k+2\iota_k\iota_k^{\mathsf T})$ of
  \eqref{eq:heights-interior-family}, in $n=2k$ variables, is computable in
  time polynomial in $k$ and has total binary length polynomial in $k$. It
  satisfies $\nabla^2h_k\succeq3I$, $h_k>0$ on $\R^n$, and
  \[
   0<m_k:=\min_{\R^n}h_k<M_k^{-2^{k+1}},\qquad M_k=1000^{k+3}.
  \]
  It is given as $\lambda_k\sum_{j=1}^{n+1}F_{k,j}^2+u_k^2$ with rational
  quadratics $F_{k,j}$, so it has a rational positive semidefinite Gram of
  rank at most $n+2<N$ and a rational sum-of-squares expression, both of
  binary length polynomial in $k$. It also has rational positive definite
  Gram matrices.
\item \emph{Lower bound.} Every rational positive definite Gram $Q$ of $h_k$,
  with reduced upper-triangular entry denominators $b_{ij}$, satisfies
  \[
   \sum_{i\le j}\log_2b_{ij}>2^k\log_2M_k-\frac{N-1}2\log_2T_k,
   \qquad T_k=\max\{1,15(n+1)\norm{h_k}_1\},
  \]
  where $\log_2T_k$ is bounded by a polynomial in $k$. Hence every rational
  positive definite Gram of $h_k$ has total denominator length
  $\Omega(k2^k)=\Omega(n2^{n/2})$, and some entry has
  $\Omega(2^k/k^3)$ denominator bits.
\end{enumerate}
\end{theorem}

The proofs are in Appendix~\ref{app:heights-interior}. The lower bound rests
on two facts. Every positive semidefinite Gram of $h_k$ has trace at most
$T_k$, because the second-moment matrix of $z$ under the uniform measure on
$[-1,1]^n$ is at least $I/(15(n+1))$. Every positive definite Gram has
smallest eigenvalue at most $m_k=\min h_k$: at a minimizer $a$, the Rayleigh
quotient on $z(a)$ is $m_k/\norm{z(a)}^2\le m_k$, since $z(a)$ contains the
constant one. Hence $0<\det Q<M_k^{-2^{k+1}}T_k^{N-1}$. On the other hand, the
product of the squared upper-triangular denominators clears every
denominator in the determinant expansion, so a positive rational determinant
is at least the reciprocal of that product. A singular Gram has determinant
zero, which is why the short Gram of rank $n+2$ escapes the bound. The upper
bound uses the strict-set rational sampling theorem of Basu, Pollack, and
Roy \cite[Theorem~4.1.2]{BasuPollackRoy1996} to find a rational center with
$c_a>0$ and $\poly(L)6^{O(n)}$ coordinate bits, and then evaluates the
explicit formula \eqref{eq:heights-completed-gram}. It is a size bound for
some interior certificate; it does not give an algorithm polynomial in $L$,
and it is not claimed for strictly positive quartics without a supplied
strict Hessian Gram.

Parts (a) and (c) show that exponential dependence on the dimension is
necessary and sufficient for interior rational Grams in this certified
class; the constants in the exponents do not match. The two kinds of Gram
matrix in the theorem should not be confused: the supplied Hessian Gram of
$h_k$ is short and has smallest eigenvalue at least three, while every
interior polynomial Gram of $h_k$ is long and has smallest eigenvalue below
$M_k^{-2^{k+1}}$. This also explains the relation to the exact rounding
theorem of G\"artner, Magron, and Vallentin
\cite[Corollary~1.3]{GaertnerMagronVallentin2026}, which recovers an exact
rational Gram when a positive lower bound on the smallest eigenvalue of some
Gram is supplied as part of the input. For $h_k$ every such bound is smaller
than $M_k^{-2^{k+1}}$, so its encoding alone has more than
$2^{k+1}\log_2M_k$ bits, and by part~(c) every rational positive definite
Gram that any method prints has $\Omega(k2^k)$ bits. Compactness of Gram
spectrahedra is
known \cite[Lemma~1.5]{ChuaPlaumannSinnVinzant2017}; the trace bound above
is an explicit version in the monomial basis that the encoding argument
needs. Exponentially large coefficients are known to be necessary in some
constrained sum-of-squares proofs \cite{ODonnell2017,RaghavendraWeitz2017};
the theorem concerns ordinary Gram matrices of a single globally positive
quartic. Whether some family in this certified class forces long
encodings of \emph{every} rational positive semidefinite Gram remains open;
part (b) shows that the family $h_k$ does not.

\subsection{Shared-circuit witnesses and interior Grams}
\label{sec:heights-circuit}

For two of the objects above, the expanded lower bounds disappear when
outputs may be shared rational circuits: strictly feasible points when the
minimum is negative, and interior Gram matrices when it is positive. The
Newton circuits of Section~\ref{sec:upper} approximate the
optimizer to doubly exponential accuracy with polynomially many gates, and the
separation bound tells how much accuracy suffices for a given polynomial
observable. With the observable $f$ itself this gives a circuit point in
$\{f<0\}$ whenever $\min f<0$
(Proposition~\ref{prop:upper-circuit-witness}); for the families $g_k$ of
Theorem~\ref{thm:heights-witness} such a point has polynomial circuit size,
while its expanded coordinates need more than $k2^k$ bits. With a different
observable the same circuit point yields interior Gram matrices, without
deciding any sign.

\begin{theorem}[Shared-circuit interior Grams]
\label{thm:circuit-gram}
There is a deterministic polynomial-time algorithm that receives a pair
$(f,A)$, where $f\in\Q[X_1,\ldots,X_n]$ has degree at most four and $A$ is a
rational symmetric matrix of order $n+n^2$, rejects the pair unless it is a
certified quartic, and otherwise outputs, without any sign oracle, a
symmetric $N\times N$ matrix $Q$ whose entries are computed by one shared
rational circuit of polynomial size, such that $f=z^{\mathsf T}Qz$
identically and
\[
 Q\succ0\quad\Longleftrightarrow\quad\min_{\R^n}f>0 .
\]
Every division in the circuit is by a nonzero rational, on every certified
input and regardless of the sign of $\min f$.
\end{theorem}

The proof is in Appendix~\ref{app:heights-circuit}. Put $\mu=\mu_A$ and
$h=f-\norm{\nabla f}^2/(2\mu)$, an explicit polynomial of degree at most six
with $h(p)=f(p)$. The shared Newton point $\hat x$ of
Lemma~\ref{lem:newton-circuit}, whose accuracy covers every observable up to
a stated encoding length, satisfies $|h(\hat x)-h(p)|\le g/8$ for a gap $g>0$
such that $h(p)\ne0$ implies $|h(p)|\ge g$. The matrix $Q$ is the completed
Taylor Gram \eqref{eq:heights-completed-gram} at the exact rational center
$\hat x$. Its constant is $c_{\hat x}=h(\hat x)$, which is positive whenever
$\min f>0$; conversely, $Q\succ0$ forces $f\ge\lambda_{\min}(Q)>0$ because
$z$ contains the constant monomial.

Neither output is an expanded rational. For the family $h_k$ of
Theorem~\ref{thm:interior-gram}, the value of $Q$ is a rational positive
definite Gram of $h_k$, so its expanded encoding needs $\Omega(n2^{n/2})$
denominator bits, while the circuit has size polynomial in $k$. For these two
outputs the circuit model removes the expanded-size obstruction, but it moves
the cost to validation over general certified quartic inputs, as follows.
\begin{enumerate}[label=(\roman*)]
\item By Theorem~\ref{thm:circuit-gram}, $Q\succ0$ holds exactly when
  $\min f>0$, and deciding $\min f>0$ for certified quartics is
  $\PosSLP$-hard (Theorem~\ref{thm:quartic-complete}). Hence deciding
  positive definiteness of a matrix given by a shared rational circuit is
  $\PosSLP$-hard already for these outputs, and it is not in polynomial time
  unless $\PosSLP$ is.
\item For the point $\hat x$ of Proposition~\ref{prop:upper-circuit-witness},
  $f(\hat x)\ge\min f$, so $f(\hat x)<0$ holds exactly when $\min f<0$.
  Checking the sign of $f(\hat x)$ is therefore as hard as deciding
  $\min f<0$, which is $\PosSLP$-complete for certified quartics.
\item The identity $f=z^{\mathsf T}Qz$ holds by construction. We claim no
  deterministic polynomial-time procedure that verifies such an identity for
  a circuit matrix received without its construction.
\item When $\min f>0$, symmetric elimination applied to $Q$ produces
  polynomial-size circuits for a weighted rational sum of squares with
  positive weights. Writing those weights as sums of rational squares by the
  binary method takes time polynomial in their \emph{expanded} length, which
  can be exponential. We do not claim polynomial-size circuits for an
  unweighted rational sum of squares.
\end{enumerate}
The construction requires the full strict Hessian Gram on $(v,X\otimes v)$:
the positive definiteness step uses $\bar A\succeq\mu\diag(0,I_{n^2})$ to
supply every translated quadratic monomial. It needs no lower bound on the
positive minimum, which may be doubly exponentially small, as for $h_k$.

The rational objects of the circle family also have short circuits. For
$f_k$ and $\tilde f_k$ of Theorem~\ref{thm:rational-height}, $k$ complex
squarings compute the optimizer, followed for $\tilde p$ by multiplications
with powers of two. The entries of $ee^{\mathsf T}$, which is the optimal
moment matrix and, up to a positive factor, every exposing matrix, are
monomials of degree at most four in the optimizer. The maximal-rank optimal
Gram $\mathcal T_p(A)$ of Theorem~\ref{thm:heights-moment}(a) is given by the
fixed formula \eqref{eq:heights-taylor-gram} in the optimizer and the
supplied Hessian Gram $A$. Hence the optimizer, $ee^{\mathsf T}$, and
$\mathcal T_p(A)$ have circuits of size polynomial in $k$, although their
expanded encodings are long (Theorem~\ref{thm:rational-height} and
Corollary~\ref{cor:heights-gram}). None of this extends to irrational
outputs. A circuit with rational constants and rational operations has a
rational value. When the optimizer is irrational, such a circuit therefore
cannot describe the optimizer, a maximal-rank optimal Gram, or a nonzero
matrix exposing the optimal Gram face
(Theorem~\ref{thm:heights-moment}(c),(e)). These objects need a different
output model, for example circuits with root gates; this section makes no
claim about such models.

\subsection{Summary of representation costs}
\label{sec:heights-summary}

Table~\ref{tab:heights-summary} collects the results. In every row the input
is a certified quartic of total length polynomial in $k$, with $n=\Theta(k)$
variables and a rational positive definite Hessian Gram.

\begin{table}[ht]
\centering
\small
\begin{tabularx}{\linewidth}{>{\raggedright\arraybackslash}X
  >{\raggedright\arraybackslash}X >{\raggedright\arraybackslash}X}
\toprule
Object & Expanded rational encoding & Short alternative\\
\midrule
Point of a strictly feasible body ($g_k$) &
every rational point has more than $\frac n2 2^{n/2}$ bits for $n\ge4$
(Thm.~\ref{thm:heights-witness}) &
circuit witness of polynomial size (Prop.~\ref{prop:upper-circuit-witness})\\
Unique rational optimizer ($f_k$, $\tilde f_k$) &
$2^{n/2-1}\log_25$ denominator bits in one coordinate
(Thm.~\ref{thm:rational-height}) &
squaring circuit of size $O(k)$; polynomial-time approximation\\
Optimal moment matrix ($f_k$) &
an entry with denominator $5^{2^k}$ (Cor.~\ref{cor:heights-gram}) &
same circuit\\
Maximal-rank optimal Gram ($f_k$) &
entry height $\Omega(2^k/n^4)$, total length $\Omega(2^k)$
(Cor.~\ref{cor:heights-gram}) &
circuit for $\mathcal T_p(A)$ from the circuit for $p$; an optimal Gram of
rank at most $n+1$ and polynomial length\\
Exposing matrix of the optimal Gram face ($f_k$) &
entry height at least $2^{k-1}\log_25$ (Cor.~\ref{cor:heights-gram}) &
$ee^{\mathsf T}$ from the circuit for $p$\\
Positive definite Gram ($h_k$) &
total denominator length $\Omega(n2^{n/2})$; at most $\poly(L)2^{O(n)}$ for
every strictly positive certified quartic (Thm.~\ref{thm:interior-gram}) &
circuit Gram of polynomial size (Thm.~\ref{thm:circuit-gram})\\
Positive semidefinite Gram or rational SOS ($h_k$) &
polynomial length (rank at most $n+2$) & ---\\
\bottomrule
\end{tabularx}
\caption{Output size of exact objects for certified quartics.}
\label{tab:heights-summary}
\end{table}

The lower bounds limit printed rational outputs only; the limits listed at
the beginning of this section apply to every row. The circuit witness and the
circuit Gram can be constructed in polynomial time. Their validation over
general certified quartic inputs has the following complexity. Deciding whether the circuit witness
$\hat x$ satisfies $f(\hat x)<0$ is $\PosSLP$-complete, and deciding whether
the circuit Gram is positive definite is $\PosSLP$-hard
(Section~\ref{sec:heights-circuit}, items (i) and (ii)). These statements
concern the general input class, not the displayed families: since
$\min g_k<0$ and $\min h_k>0$, both predicates are always true on those
families.
